\documentclass[../DoAn.tex]{subfiles}
\begin{document}

Chương 5 đã phân tích chi tiết kết quả thực nghiệm, đánh giá các chỉ số hiệu năng đo lường thực tế và xác định năng lực phòng chống các lỗ hổng bảo mật của hệ thống phân tán theo tiêu chuẩn OWASP Top 10. Chương 6 này là chương cuối cùng của báo cáo đồ án tốt nghiệp, chịu trách nhiệm đưa ra các kết luận tổng kết toàn diện về toàn bộ quá trình nghiên cứu, thiết kế và triển khai giải pháp hệ thống điểm danh thông minh. Sự có mặt của chương này giúp nhìn nhận lại một cách tổng quan các thành tựu kỹ thuật đã đạt được so với mục tiêu ban đầu, đồng thời định hình rõ nét các không gian cải tiến và vạch ra định hướng công việc trong tương lai nhằm nâng cấp sản phẩm phần mềm. Nội dung cốt lõi của chương được chia làm hai phần chính bao gồm: (i) tổng kết các kết quả kỹ thuật nổi bật, đối chiếu sản phẩm với các nghiên cứu tương tự và đánh giá bài học kinh nghiệm rút ra trong phần 6.1, và (ii) phân tích các điểm giới hạn kỹ thuật còn tồn tại, từ đó đề xuất các hướng nghiên cứu, nâng cấp chuyên sâu cho hệ thống trong phần 6.2.

\section{Kết luận}
\label{section:6.1}

\subsection{Các kết quả kỹ thuật đã đạt được (Đã làm được)}
Đồ án đã nghiên cứu, thiết kế và cài đặt hiện thực hóa thành công một hệ thống điểm danh thông minh tích hợp nhận diện khuôn mặt và bóc tách dữ liệu chip CCCD Việt Nam, đáp ứng đầy đủ và toàn diện các mục tiêu khoa học, kỹ thuật đã đề ra trong Chương 1. Hệ thống bao gồm năm thành phần hạt nhân hoạt động phối hợp chặt chẽ, nhịp nhàng:
\begin{itemize}
    \item \textbf{Ứng dụng di động Flutter (Trạm đầu cuối):} Thực hiện nghiệp vụ bóc tách NFC theo giao thức BAC/PACE, đối soát lớp 1 bằng MobileFaceNet chạy offline trực tiếp trên thiết bị thông qua ONNX Runtime và hỗ trợ hiển thị đồ họa banner nhắc nhở vi phạm chuỗi trực quan ngay tại hiện trường.
    \item \textbf{Lõi máy chủ Backend Spring Boot:} Điều phối toàn bộ luồng xác thực lớp 2 nâng cao, thực thi cơ chế \textbf{An ninh mềm đối xứng hai chiều (Biometric-Driven Soft Mode)} dựa theo mốc giờ \texttt{end\_time} cá nhân, tự động quản lý đồng bộ ma trận phân quyền theo cặp thiết bị và xử lý các logic chấm công hành chính.
    \item \textbf{Dịch vụ tính toán Python Flask:} Tiếp nhận request payload, bóc tách chuỗi hình ảnh và xử lý các tác vụ so sánh độ tương đồng hình học không gian Cosine bằng mô hình học sâu chuyên sâu ArcFace với độ ổn định cao.
    \item \textbf{Bo mạch vi điều khiển ESP32:} Kết nối kéo dài qua giao thức TCP ổn định thông qua mạng đám mây HiveMQ Cloud cổng bảo mật 8883, đóng vai trò cơ cấu chấp hành nhận lệnh từ máy chủ để đóng ngắt rơ-le khóa cửa vật lý một cách an toàn.
    \item \textbf{Giao diện Web React:} Cung cấp hệ thống quản trị hành chính \textbf{vận hành tại cổng mạng 3002}, cho phép cấu hình ma trận phân quyền, xuất báo cáo Excel và hiển thị Web Dashboard giám sát dòng người ra vào thời gian thực, tự động phân loại trực quan các Badge sự cố màu cam (\texttt{MISSING\_IN\_AUTO\_OPEN} và \texttt{ANTI\_PASSBACK\_IN}).
\end{itemize}

Về mặt học thuật và kỹ thuật, đồ án đã vận dụng thành công nhiều giải pháp công nghệ hiện đại trong một cấu trúc kiến trúc phân tán mạch lạc, tối ưu: ứng dụng mô hình học sâu InsightFace đạt độ chính xác cao; giao thức bảo mật ICAO Doc 9303 cho phân hệ đọc chip thẻ vật lý; giao thức mạng nhẹ MQTT over SSL/TLS cổng 8883 cho phân hệ điều khiển phần cứng; giải pháp mã hóa phi trạng thái JWT phối hợp bộ lọc Spring Security và kiểm soát đặc quyền \texttt{@PreAuthorize} (RBAC năm cấp bao gồm \texttt{SUPER\_ADMIN}, \texttt{ADMIN}, \texttt{HR\_MANAGER}, \texttt{OPERATOR}, \texttt{VIEWER}) cho hành lang bảo mật REST API; và đường ống song công WebSocket STOMP cho phân hệ Web Dashboard cổng 3002.

Kết quả đo lường thực nghiệm khẳng định hệ thống đạt được sự cân bằng tối ưu giữa ba tiêu chí: độ chính xác sinh trắc học cao (94.5\% TPR, 1.2\% FPR với ngưỡng quyết định 0.65), tốc độ xử lý phản hồi thời gian thực ấn tượng (đạt trung bình 3.3 giây cho toàn bộ chu trình đầu cuối chuẩn và 3.8 giây cho luồng tích hợp Soft Mode) và tính an toàn thông tin vững chắc bám sát tiêu chuẩn OWASP Top 10.

\subsection{Hạn chế của hệ thống thực nghiệm (Chưa làm được)}
Mặc dù đạt được các chỉ số hiệu năng ấn tượng, hệ thống vẫn tồn tại một số hạn chế kỹ thuật chưa thể giải quyết triệt để trong phạm vi thời gian thực hiện đồ án:
\begin{itemize}
    \item \textbf{Sự phụ thuộc hạ tầng phần cứng của trạm đầu cuối:} Phân hệ bóc tách dữ liệu DG2 bắt buộc thiết bị di động Android của người dùng phải tích hợp phần cứng đầu đọc NFC ổn định và camera trước độ phân giải cao, gây khó khăn cho việc triển khai đồng bộ trên các dòng máy phân khúc thấp.
    \item \textbf{Chưa thực hiện xác thực Passive Authentication chuyên sâu:} Hệ thống mới chỉ dừng lại ở bước bắt tay thiết lập kênh truyền bảo mật BAC để đọc dữ liệu từ chip, chưa xây dựng phân hệ kiểm tra chuỗi chứng chỉ số SOD (Passive Authentication) để đối soát chữ ký điện tử của Cơ quan cấp phát với Root CA quốc gia.
    \item \textbf{Khả năng chống tấn công giả mạo sinh trắc học vật lý còn sơ khai:} Bộ lọc camera selfie mới chỉ nhận diện cấu trúc khuôn mặt người sống, chưa tích hợp mô hình phát hiện thực thể sống (Liveness Detection) chuyên sâu để ngăn chặn các hình thức tấn công dùng ảnh in hoặc video phát lại từ màn hình.
\end{itemize}

\subsection{So sánh, đối chiếu kết quả với các nghiên cứu tương tự}
Để làm rõ giá trị khoa học và thực tiễn của đề tài, sản phẩm đồ án được tiến hành đối chiếu toàn diện với các giải pháp điểm danh thông dụng hiện nay trên thị trường (như máy chấm công vân tay/khuôn mặt truyền thống độc lập) và các giải pháp định danh camera AI tập trung.

\begin{itemize}
    \item \textbf{So với giải pháp máy chấm công độc lập truyền thống (Vân tay/Thẻ từ nội bộ):} Các hệ thống truyền thống có chi phí đầu tư thiết bị rất cao tại từng cửa, dễ bị gian lận bằng hình thức quẹt thẻ hộ, và tốn kém chi phí logistics khi cấp phát/thu hồi thẻ nội bộ. Hệ thống của đề tài tận dụng trực tiếp hạ tầng định danh quốc gia có sẵn (thẻ CCCD gắn chip), giúp tổ chức triệt tiêu hoàn toàn 100\% chi phí mua sắm vật tư thẻ, đồng thời ngăn chặn tuyệt đối hiện tượng chấm công hộ nhờ cơ chế đối soát sinh trắc học bắt buộc đi kèm.
    \item \textbf{So với giải pháp nhận diện Camera AI tập trung (Centralized Face Recognition Cloud):} Mô hình camera AI tập trung yêu cầu máy chủ phải lưu trữ một kho dữ liệu ảnh chân dung gốc của toàn bộ nhân sự để làm mẫu đối sánh. Mô hình này đối mặt với nguy cơ rò rỉ dữ liệu cá nhân rất lớn (Single Point of Failure) và giảm độ chính xác theo thời gian do hiện tượng lão hóa khuôn mặt (Aging effect). Giải pháp kiến trúc 2 lớp tự trị của đồ án chuyển dịch hoàn toàn sang cơ chế kiểm tra chéo động (Dynamic Cross-Verification). Bằng cách bóc tách trực tiếp ảnh tham chiếu chuẩn từ chip CCCD tại thời điểm quét và truyền trong request payload để đối soát chéo với ảnh selfie, hệ thống loại bỏ hoàn toàn việc xây dựng kho lưu trữ ảnh tập trung, giải quyết triệt để bài toán bảo mật thông tin và rủi ro pháp lý về quyền riêng tư.
\end{itemize}

\subsection{Bài học kinh nghiệm rút ra}
Quá trình nghiên cứu và phát triển một hệ thống phân tán đa nền tảng, tích hợp sâu giữa phần cứng nhúng, dịch vụ học sâu và các framework phần mềm hiện đại đã mang lại cho sinh viên những bài học kinh nghiệm chuyên môn vô cùng quý báu:
\begin{enumerate}
    \item \textbf{Tư duy thiết kế hệ thống bất đồng bộ và tự phục hồi chuỗi:} Việc xử lý luồng truyền thông đa thiết bị đòi hỏi phải từ bỏ tư duy đồng bộ truyền thống. Giải pháp tối ưu hóa hồ chứa luồng (Thread Pool) phối hợp cấu trúc Async Executor của Spring Boot và đường ống truyền dữ liệu WebSocket STOMP là chìa khóa để triệt tiêu hiện tượng nghẽn mạch cục bộ. Việc tích hợp thành công Soft Mode chứng minh tư duy thiết kế hệ thống hướng người dùng, cân bằng hoàn hảo giữa an ninh chặt chẽ và lưu thông thực địa.
    \item \textbf{Tầm quan trọng của thiết kế cảnh báo trực quan (Visual Alert):} Để bám sát môi trường vận hành văn phòng chuyên nghiệp, đồ án rút ra bài học lớn về việc chuyển dịch các cơ chế còi báo động vật lý thô sơ sang hạ tầng cảnh báo đồ họa trực quan. Sự phối hợp song công giữa banner thông báo của Android Client và hệ thống định danh màu sắc Badge cam trên Web cổng 3002 tối ưu hóa năng lực giám sát từ xa cho OPERATOR một cách tinh tế.
    \item \textbf{Kỹ năng tối ưu hóa tài nguyên phần cứng mã nguồn:} Khi triển khai các mô hình học sâu lên môi trường biên (Edge Computing) như thiết bị di động, việc chuyển đổi trọng số sang định dạng rút gọn mã nguồn mở và chạy qua thư viện ONNX Runtime giúp giảm thiểu đáng kể dung lượng RAM chiếm dụng và tăng tốc độ xử lý toán học mà không làm suy giảm quá nhiều độ chính xác của mô hình.
\end{enumerate}

\section{Hướng phát triển tương lai}
\label{section:6.2}

Để khắc phục triệt để các giới hạn kỹ thuật còn tồn tại, đồng thời nâng cao giá trị thực tiễn và năng lực cạnh tranh công nghệ của sản phẩm, hướng đi tiếp theo của đề tài tập trung vào hai nhóm nhiệm vụ cốt lõi:

\subsection{Hoàn thiện và nâng cấp các chức năng hiện tại}
\begin{itemize}
    \item \textbf{Hiện thực hóa toàn diện thuật toán Passive Authentication:} Nghiên cứu cài đặt quy trình bóc tách và xác chuẩn phân vùng tệp tin SOD theo đúng đặc tả kỹ thuật ICAO Doc 9303. Phân hệ này sẽ nạp danh sách chứng chỉ số công khai Root CA của Bộ Công an Việt Nam để kiểm tra chữ ký điện tử, giúp máy chủ phát hiện và chặn đứng hoàn toàn các nguy cơ sử dụng chip thẻ giả mạo hoặc tệp tin mã nhị phân đã bị sửa đổi.
    \item \textbf{Tích hợp giải pháp chống giả mạo sinh trắc học chuyên sâu (Liveness Detection):} Nghiên cứu tích hợp các thuật toán phát hiện thực thể sống dựa trên việc phân tích cấu trúc kết cấu da, chuyển động vi mô của cơ mắt hoặc phản xạ ánh sáng tần số quét nhằm nâng cao năng lực an ninh thực địa, chặn đứng các hành vi tấn công giả mạo dạng trình chiếu sử dụng ảnh in hoặc video phát lại.
    \item \textbf{Tinh chỉnh mô hình AI trên dữ liệu đặc trưng nội địa:} Tổ chức thu thập tập mẫu khuôn mặt người Việt Nam và thực hiện kỹ thuật tinh chỉnh mô hình (Fine-tuning) trên mô hình nền ArcFace nhằm tối ưu hóa năng lực bóc tách vector đặc trưng cho các đặc điểm hình học khuôn mặt đặc thù của người nội địa, đẩy chỉ số TPR lên mức trên 97\%.
\end{itemize}

\subsection{Nghiên cứu các hướng mở rộng, nâng cấp hệ thống chuyên sâu}
\begin{itemize}
    \item \textbf{Nâng cấp hạ tầng triển khai lên kiến trúc phân tán quy mô lớn:} Để chuẩn bị cho kịch bản đưa vào vận hành thực tế tại các tập đoàn có hàng vạn nhân sự và hàng trăm trạm cửa vật lý chạy đồng thời, toàn bộ hạ tầng cần được Container hóa bằng Docker, điều phối tự động bằng Kubernetes hỗ trợ cơ chế tự động co giãn (Horizontal Pod Autoscaling). Hệ thống tận dụng tính năng \textbf{Auto-create Permission} theo cặp thiết bị đối xứng dựa theo trường \texttt{location\_name} của bảng \texttt{devices} để làm bệ phóng triển khai hành chính quy mô lớn tự trị.
    \item \textbf{Hỗ trợ luồng điểm danh nhanh kết hợp:} Bổ sung module chức năng hỗ trợ luồng điểm danh nhanh chỉ bằng sinh trắc học khuôn mặt (1:N) tại các chu kỳ sau khi nhân viên đã hoàn thành bước xác thực thẻ CCCD lần đầu, giúp tăng tốc độ lưu thông dòng người vào các khung giờ cao điểm.
    \item \textbf{Liên kết hệ thống quản trị nhân sự tập trung:} Thiết lập hệ thống giao diện lập trình để liên kết API đồng bộ dữ liệu chấm công trực tiếp với các phần mềm quản trị nguồn nhân lực (HRM) phổ biến tại Việt Nam như VnResource hoặc MISA HRM.
    \item \textbf{Phân tích xu hướng bằng Học máy (Machine Learning):} Tích hợp các thuật toán học máy ở lớp tầng dữ liệu để phân tích, dự báo các chỉ số xu hướng hành vi đi muộn, vắng mặt của nhân sự theo mùa hoặc theo phòng ban, cung cấp góc nhìn toàn diện cho nhà quản trị.
\end{itemize}

\end{document}